\documentclass[12pt,oneside]{book}

% ===================== METADATOS =====================
\newcommand{\universidad}{Universidad de Buenos Aires (UBA)}
\newcommand{\facultad}{Facultad de Derecho}
\newcommand{\materia}{Compraventa Inmobiliaria}
\newcommand{\titulo}{Apuntes de Compraventa Inmobiliaria}
\newcommand{\autor}{Isaac Edgar Camacho Ocampo}
\newcommand{\anio}{2026}

% ===================== PAQUETES =====================
\usepackage[utf8]{inputenc}
\usepackage[T1]{fontenc}
\usepackage[spanish,es-nodecimaldot]{babel}
\usepackage[top=2.5cm,bottom=2.5cm,left=3cm,right=2.5cm,headheight=15pt]{geometry}
\usepackage{amsmath}
\usepackage{amssymb}
\usepackage{mathtools}
\usepackage{graphicx}
\usepackage{float}
\usepackage{caption}
\usepackage{subcaption}
\usepackage{booktabs}
\usepackage{array}
\usepackage{longtable}
\usepackage{tabularx}
\usepackage{enumitem}
\usepackage{xcolor}
\usepackage[most]{tcolorbox}
\usepackage{fancyhdr}
\usepackage{ragged2e}
\usepackage[
    colorlinks=true,
    linkcolor=blue!50!black,
    citecolor=green!40!black,
    urlcolor=blue!60!black,
    pdftitle={\titulo},
    pdfauthor={\autor},
    pdfsubject={Apuntes universitarios},
    bookmarks=true
]{hyperref}

% ===================== CONFIGURACIÓN =====================
\graphicspath{{./img/}{./graficos/}{./figuras/}}
\setcounter{tocdepth}{2}
\setlength{\parindent}{0pt}
\setlength{\parskip}{0.5em}
\setlist[itemize]{topsep=0.3em,itemsep=0.2em}
\setlist[enumerate]{topsep=0.3em,itemsep=0.2em}
\clubpenalty=10000
\widowpenalty=10000

% ===================== COMANDOS =====================
\newcommand{\abs}[1]{\left|#1\right|}
\newcommand{\norm}[1]{\left\lVert#1\right\rVert}
\newcommand{\vect}[1]{\mathbf{#1}}
\newcommand{\deriv}[2]{\frac{d#1}{d#2}}
\newcommand{\pderiv}[2]{\frac{\partial#1}{\partial#2}}

% ===================== CAJAS =====================
\newtcolorbox{definition}[1]{enhanced,breakable,title={Definición: #1},fonttitle=\bfseries,colback=blue!3,colframe=blue!50!black,boxrule=0.8pt,arc=2mm}
\newtcolorbox{important}{enhanced,breakable,title={Importante},fonttitle=\bfseries,colback=yellow!8,colframe=orange!70!black,boxrule=0.8pt,arc=2mm}
\newtcolorbox{example}[1]{enhanced,breakable,title={Ejemplo: #1},fonttitle=\bfseries,colback=green!4,colframe=green!40!black,boxrule=0.8pt,arc=2mm}
\newtcolorbox{formula}[1]{enhanced,breakable,title={Fórmula / Regla: #1},fonttitle=\bfseries,colback=purple!4,colframe=purple!50!black,boxrule=0.8pt,arc=2mm}
\newtcolorbox{exercise}[1]{enhanced,breakable,title={Ejercicio: #1},fonttitle=\bfseries,colback=gray!5,colframe=gray!60!black,boxrule=0.8pt,arc=2mm}
\newtcolorbox{note}{enhanced,breakable,title={Nota},fonttitle=\bfseries,colback=white,colframe=gray!50,boxrule=0.6pt,arc=2mm}

% ===================== ENCABEZADOS =====================
\pagestyle{fancy}
\fancyhf{}
\fancyhead[L]{\nouppercase{\leftmark}}
\fancyhead[R]{\materia}
\fancyfoot[C]{\thepage}
\renewcommand{\headrulewidth}{0.4pt}
\renewcommand{\footrulewidth}{0pt}
\fancypagestyle{plain}{\fancyhf{}\fancyfoot[C]{\thepage}\renewcommand{\headrulewidth}{0pt}}

% ===================== DOCUMENTO =====================
\begin{document}

\begin{titlepage}
\thispagestyle{empty}
\begin{center}
\vspace*{1cm}
{\large \textsc{\universidad}}\\[0.2cm]
{\large \textsc{\facultad}}\\[3cm]
{\Huge \textbf{\materia}}\\[1cm]
{\LARGE \titulo}\\[0.8cm]
\textit{Comprender los conceptos es el primer paso para convertir el estudio en conocimiento.}\\[1.2cm]
\rule{0.70\textwidth}{0.5pt}\\[0.5cm]
{\large Apuntes de estudio}\\[0.5cm]
\rule{0.70\textwidth}{0.5pt}
\vfill
{\large \autor}\\[0.3cm]
{\large Octubre de \anio}
\end{center}
\vfill
\begin{flushleft}
\textit{Este es un modesto aporte a los estudiantes de ``Compraventa Inmobiliaria'' no pretende sustituir el material oficial de la materia.}
\end{flushleft}
\end{titlepage}

\tableofcontents

% =====================================================
\chapter{El objeto del contrato de compraventa}
% =====================================================

Esta unidad estudia el objeto del contrato de compraventa de inmuebles: qué cosas pueden venderse, qué ocurre cuando el inmueble tiene derechos de terceros o medidas judiciales y cómo se identifica el inmueble que se vende. Se parte de la estructura de los contratos en general para luego llevarla a los inmuebles en particular.

Los artículos de referencia son los arts.\ 1129, 16, 279, 1003, 1004 y 1009 del Código Civil y Comercial.

\section{Cosa vendida y marco normativo}

\subsection{La cosa vendida y el objeto del acto jurídico}

El art.\ 1129 establece que pueden venderse todas las cosas que pueden ser objeto de los contratos. Por ello debe acudirse a la parte general para determinar qué puede ser objeto de un contrato.

\begin{definition}{Art. 1129 CCyCN --- Cosa vendida}
Pueden venderse todas las cosas que pueden ser objeto de los contratos.
\end{definition}

\begin{definition}{Art. 279 CCyCN --- Objeto}
El objeto del acto jurídico no debe ser un hecho imposible o prohibido por la ley, contrario a la moral, a las buenas costumbres, al orden público o lesivo de los derechos ajenos o de la dignidad humana. Tampoco puede ser un bien que por un motivo especial se haya prohibido que lo sea.
\end{definition}

\subsection{El objeto del contrato y el concepto de «cosa»}

\begin{definition}{Art. 1003 CCyCN --- Disposiciones generales}
Se aplican al objeto del contrato las disposiciones de la Sección 1ª, Capítulo 5, Título IV del Libro Primero de este Código. Debe ser lícito, posible, determinado o determinable, susceptible de valoración económica y corresponder a un interés de las partes, aun cuando éste no sea patrimonial.
\end{definition}

\begin{definition}{Art. 16 CCyCN --- Bienes y cosas}
Los derechos referidos en el primer párrafo del artículo 15 pueden recaer sobre bienes susceptibles de valoración económica. Los bienes materiales se llaman cosas en este Código. Las disposiciones referentes a las cosas son aplicables a la energía y a las fuerzas naturales susceptibles de ser puestas al servicio del hombre.
\end{definition}

Objeto de los contratos son todas las cosas que puedan ser vendidas. Las cosas son bienes susceptibles de valoración económica o pecuniaria. Para que exista contrato, y particularmente contrato de compraventa como contrato oneroso, el objeto debe ser susceptible de apreciación pecuniaria. En el caso de los inmuebles no hay duda alguna: forman parte del patrimonio de las personas y, en principio, cumplen con ese texto.

\section{Objetos prohibidos y dominio público}

\begin{definition}{Art. 1004 CCyCN --- Objetos prohibidos}
No pueden ser objeto de los contratos los hechos que son imposibles o están prohibidos por las leyes, son contrarios a la moral, al orden público, a la dignidad de la persona humana, o lesionan derechos ajenos; ni los bienes que por un motivo especial se prohíbe que lo sean. Cuando tengan por objeto derechos sobre el cuerpo humano se aplican los artículos 17 y 56.
\end{definition}

Conforme al art.\ 1004, la prestación no debe ser de cumplimiento imposible, ni referirse a una cosa prohibida por las leyes, ni ser contraria a la moral y a las buenas costumbres. Ciertos inmuebles, como los \textbf{inmuebles de dominio público del Estado}, no pueden ser objeto de un contrato celebrado entre las partes, porque estaría prohibido por la ley.

Se analizan luego los inmuebles de \textbf{dominio privado} sobre los que recae alguna medida cautelar, algún derecho real de disfrute o algún derecho real accesorio de garantía, como la hipoteca. Corresponde determinar si los inmuebles sobre los que recaen estos gravámenes y medidas pueden ser objeto o no de un contrato de compraventa.

\section{Requisitos del objeto}

\subsection{El objeto debe estar determinado o ser determinable}

Se exige un criterio mínimo de determinación: la cosa, al momento de celebrarse el contrato, debe estar determinada o ser determinable.

\begin{itemize}
\item \textbf{Cosa determinada:} se puede individualizar la cosa, específicamente su especie y su género.
\item \textbf{Cosa determinable:} se establecen pautas o criterios suficientes para individualizar la cosa a futuro, con posterioridad.
\end{itemize}

\begin{example}{Inmueble en pozo}
Es frecuente la compra de inmuebles de pozo, en los que puede existir una finca o un terreno que todavía no se afectó al régimen de propiedad horizontal. Todavía no habrá nacido la unidad funcional cinco que se quiere comprar. En ese caso se dejan pautas o criterios suficientes para establecer que, a futuro, se comprará efectivamente esa unidad. Por ejemplo, la unidad que se encuentra en el piso cuarto, con frente a la avenida o en contrafrente. Que luego cambie o no el número es otra cuestión, pero se dejan los lineamientos básicos para que, cuando la cosa exista verdaderamente, pueda determinarse e identificarse correctamente.
\end{example}

\textbf{¿Por qué es importante la determinación?}
\begin{enumerate}
\item Para que no quepa duda sobre qué se obliga a transmitir el vendedor y qué es lo que compra el comprador a cambio de la suma de dinero que entrega. En el ámbito inmobiliario, muchas veces se trata de los ahorros de toda una vida, por lo que la individualización es muy importante.
\item Por cuestiones de seguridad jurídica relacionadas con la causa: en los contratos onerosos, como la compraventa, debe haber igualdad y equilibrio en las prestaciones. No se puede comprar un inmueble por 10.000 pesos sin que las consecuencias sean diferentes; debe haber correlatividad en las prestaciones. Las características del precio se verán más adelante.
\end{enumerate}

\subsection{La cosa debe tener existencia real o posible}

La existencia es \textbf{real} cuando efectivamente está la cosa, y \textbf{posible} cuando hay al menos un grado de certeza de que la cosa pueda llegar a existir en algún momento.

\begin{example}{Departamentos en pozo}
Quien va a comprar un inmueble en pozo suele hablar con la desarrolladora, que muestra de forma digital los planos, el AutoCAD o el render del edificio. Según la instancia en que se encuentre, puede haberse ingresado ya con el hormigón, lo que da ciertos inicios y cierto grado de certeza de que el inmueble llegará a existir.
\end{example}

\section{Utilidad del tema}

Comprar una vivienda suele ser la operación más grande de la vida de una persona. Saber qué preguntar frente a un certificado con hipoteca, embargo o usufructo protege el patrimonio. Si se compra con un embargo vigente se asume el riesgo y no se puede reclamar al vendedor; si se asume una deuda hipotecaria se responde con todo el patrimonio, y si solo se la reconoce se responde únicamente con el inmueble.

Aplicaciones profesionales señaladas en el apunte:
\begin{itemize}
\item Leer un certificado o informe de dominio y decidir qué se puede escriturar, qué debe hacer el comprador (notificar, reconocer o asumir) y qué trámites pedir antes de la firma.
\item Redactar y controlar cláusulas de escrituras y boletos: reconocimiento de servidumbre, «reconoce el embargo al solo efecto registral», declaración jurada del vendedor, descripción de la unidad funcional.
\item Asesorar con criterio de riesgo: distinguir las medidas que impiden vender de las que solo eliminan la buena fe, y evitar que el Registro devuelva o observe una escritura.
\end{itemize}

\section*{Cuadro Cornell: el objeto del contrato}
\addcontentsline{toc}{section}{Cuadro Cornell: el objeto del contrato}

\begin{longtable}{>{\raggedright\arraybackslash}p{0.27\textwidth}>{\raggedright\arraybackslash}p{0.65\textwidth}}
\toprule
\textbf{Preguntas / palabras clave} & \textbf{Notas / respuestas} \\
\midrule
\endfirsthead
\toprule
\textbf{Preguntas / palabras clave} & \textbf{Notas / respuestas} \\
\midrule
\endhead
\textbf{¿Qué cosas pueden ser objeto de una compraventa? (arts.\ 1129, 279, 1003, 16)} &
Pueden venderse todas las cosas que pueden ser objeto de los contratos (art.\ 1129). El objeto debe ser lícito, posible, determinado o determinable, susceptible de valoración económica y corresponder a un interés de las partes (art.\ 1003); no puede ser un hecho imposible o prohibido, ni un bien que se prohibió que lo sea (arts.\ 279 y 1004). Las cosas son los bienes materiales susceptibles de valoración económica (art.\ 16). Los inmuebles cumplen estos requisitos, salvo los del dominio público del Estado. \\[0.6em]
\textbf{¿Qué significan «determinado» y «determinable»?} &
Determinada: se puede individualizar la cosa (su especie, su género). Determinable: se fijan pautas o criterios suficientes para individualizarla más adelante. Ejemplo: una unidad funcional que aún no nació porque el terreno no está afectado a propiedad horizontal; se detalla, por ejemplo, el piso y el frente o contrafrente. \\[0.6em]
\textbf{¿Por qué importan la determinación y la existencia real o posible?} &
Para que no haya duda sobre qué transmite el vendedor y qué compra el comprador; por seguridad jurídica; y por la correlatividad entre las prestaciones en los contratos onerosos. La cosa debe existir (real) o haber certeza de que llegará a existir (posible): inmuebles en pozo con planos, render u obra iniciada. \\
\midrule
\multicolumn{2}{p{0.94\textwidth}}{\textbf{Resumen:} un inmueble puede ser objeto de compraventa si es lícito, posible, determinado o determinable, valuable económicamente y de existencia real o posible; quedan excluidos los inmuebles del dominio público del Estado.} \\
\bottomrule
\end{longtable}

% =====================================================
\chapter{Inmuebles litigiosos, gravados o con medidas cautelares: derechos de terceros}
% =====================================================

\section{El artículo 1009 y la regla general}

\begin{definition}{Art. 1009 CCyCN --- Bienes litigiosos, gravados, o sujetos a medidas cautelares}
Los bienes litigiosos, gravados, o sujetos a medidas cautelares, pueden ser objeto de los contratos, sin perjuicio de los derechos de terceros. Quien de mala fe contrata sobre esos bienes como si estuviesen libres debe reparar los daños causados a la otra parte si ésta ha obrado de buena fe.
\end{definition}

«Gravado» significa que la constitución del gravamen es voluntaria; la medida cautelar, en cambio, es obligatoria y la impone un juez.

\begin{important}
Regla general del art.\ 1009: los bienes litigiosos, gravados o sujetos a medidas cautelares pueden ser objeto de los contratos; se pueden vender. Siempre debe cuidarse de no alterar ni perjudicar el derecho de los terceros. Pueden ser terceros el usufructuario (derechos reales de disfrute), el acreedor hipotecario (hipoteca) y el acreedor (medida cautelar).
\end{important}

\section{Las tres situaciones posibles}

Sobre un inmueble pueden darse tres situaciones:
\begin{enumerate}
\item Que recaigan \textbf{derechos reales de disfrute}.
\item Que recaigan \textbf{derechos reales de garantía}, otro gravamen de constitución voluntaria.
\item Que recaigan \textbf{medidas cautelares}. Se distinguen las medidas cautelares reales, que recaen sobre la cosa, de las personales, que recaen sobre la persona (la inhibición recae sobre la persona).
\end{enumerate}

La mejor situación es que el certificado del inmueble figure completamente limpio: sin medidas cautelares, sin derechos reales de disfrute y sin derechos reales de garantía. Es una situación ideal que en la práctica no sucede; siempre debe resolverse alguna cuestión. Se analiza entonces si los inmuebles con estas situaciones pueden venderse y, de ser así, qué hay que hacer.

\section{Derechos reales de disfrute}

\subsection{Usufructo, uso y habitación: la venta de la nuda propiedad}

Las reglas del usufructo se aplican luego al uso y a la habitación. Si el certificado figura un usufructo, el inmueble puede venderse, pero se vende la \textbf{nuda propiedad}, porque probablemente exista un dominio desmembrado.

\begin{example}{Nudo propietario y usufructuario}
El nudo propietario puede vender su nuda propiedad sin requerir la conformidad del usufructuario, pero debe notificarlo. El usufructuario debe reconocer el \textit{animus domini} en otra persona; seguirá teniendo su derecho real de usufructo, que no se altera. Como el dominio está desmembrado, el nudo propietario puede disponer de su nuda propiedad.
\end{example}

Respecto del modo (el derecho real nace con título suficiente y modo suficiente): el usufructuario ocupa el inmueble y tiene la disposición material y jurídica de la cosa. Ambos son poseedores a nombre de distinto título: el usufructuario por su usufructo y el nudo propietario por su nuda propiedad.

Si el usufructuario quisiera transmitir su usufructo, tampoco hay problema; el nuevo usufructuario debe reconocer al nudo propietario como tal y otorgar una garantía suficiente.

\begin{important}
En ambos casos (venta de la nuda propiedad o del usufructo) siempre hay que \textbf{notificar} a la otra parte. \textbf{No se requiere conformidad}, a diferencia de lo que ocurre en los derechos de garantía. El nudo propietario puede disponer libremente de su derecho, sin permiso del usufructuario, titular de un derecho de menor jerarquía; la notificación sirve para que el usufructuario sepa que existe otra persona con un derecho superior. En ambas situaciones cada uno puede disponer de su derecho real a título gratuito o a título oneroso.
\end{important}

\subsection{Fallecimiento del usufructuario o del nudo propietario}

Si fallece el usufructuario y el usufructo es gratuito y vitalicio, el nudo propietario adquiere el dominio pleno, a menos que exista el derecho de acrecer de otro. Por ejemplo, si en lugar de un solo usufructuario son dos (un matrimonio), al fallecer uno acrece el otro; todo ello debe estar previsto. En el caso del nudo propietario, se produce una sucesión. Siempre se notifica a la otra parte. Para el uso y la habitación rige la misma lógica.

\subsection{Servidumbres}

Existen servidumbres de paso y servidumbres administrativas. En la práctica sucede cuando se vende una unidad complementaria o una unidad funcional destinada a cochera, a la que para acceder hay que pasar por la de otra persona; o, por ejemplo, el baño del encargado.

Del informe o certificado registral debe surgir la servidumbre, y lo único que hace el comprador es \textbf{reconocerla}: acepta y reconoce la servidumbre que surge del certificado tal o cual. En la escritura, al consignar los certificados registrales, se indica que de tal certificado surge una servidumbre que el comprador manifiesta reconocer y aceptar en ese acto. Al inscribir la escritura en el Registro de la Propiedad Inmueble se agrega como acto el reconocimiento de una servidumbre, porque la servidumbre también es un derecho real de disfrute. No implica mayores complicaciones.

\section{Derechos reales de garantía: la hipoteca}

La hipoteca es uno de los derechos reales de garantía más comunes.

\subsection{Formas de estructurar la hipoteca}

\begin{itemize}
\item \textbf{Con un acreedor externo (un tercero):} por ejemplo, un crédito bancario para la primera vivienda, en el que el acreedor hipotecario es el banco. El negocio jurídico causal es un contrato de mutuo, un préstamo de dinero.
\item \textbf{Compraventa con hipoteca por saldo de precio:} el vendedor ofrece un inmueble a \$100.000, el comprador tiene \$50.000 y se pacta que los \$50.000 restantes se entreguen en cinco cuotas de \$10.000. El vendedor es entonces el acreedor hipotecario.
\end{itemize}

\subsection{Las tres alternativas frente al inmueble hipotecado}

\textbf{Primera alternativa: cancelar la hipoteca.} El vendedor cancela su deuda con el acreedor al vender. El comprador recibe el inmueble libre de cualquier gravamen. Es la mejor situación: acto uno, cancelo; acto dos, vendo libre. Si la deuda ya había sido abonada, la información del registro está desactualizada porque informa una hipoteca que garantiza una deuda que ya no existe.

\textbf{Segunda alternativa: reconocer y aceptar la hipoteca.} El adquirente reconoce la existencia de la hipoteca, pero quien sigue pagando es el deudor original, que no se libera. Ante una eventual ejecución hipotecaria, el comprador responde por el límite del gravamen o del valor del bien, no con todo su patrimonio.

\textbf{Tercera alternativa: asumir la deuda.} El comprador ocupa el lugar del deudor inicial. Como deudor, todo su patrimonio es prenda común del acreedor.

\begin{example}{Reconocer o asumir}
Si el inmueble está hipotecado por 15 millones y se compró a 20 millones, corresponde responder por 15. Si el inmueble vale 100 y la deuda total es 120: quien \textit{reconoce} la hipoteca responde solo con el inmueble y no debe 20, porque no asumió la deuda; quien \textit{asume} la deuda responde por 120, con todo su patrimonio.
\end{example}

\begin{important}
\textbf{¿Cuándo tiene que venir el acreedor?} Solo en caso de \textbf{delegación pasiva de deuda}: el acreedor exonera al deudor original, lo libera, y el comprador pasa a ser el único deudor. Para eso debe asistir a la escritura y prestar su conformidad, porque el patrimonio es prenda común de todos los acreedores y al acreedor no le da lo mismo cobrar de una persona que de otra. Si no asiste y el comprador asumió la deuda, el acreedor tendrá dos patrimonios para ejecutar: el del deudor originario, que contrajo la deuda, y el del comprador, que la asumió. La conformidad se presta solo para exonerar.
\end{important}

\subsection{Consulta sobre la cancelación simultánea}

Un alumno consultó si en el primer supuesto (venta con cancelación de hipoteca) puede hacerse por tracto abreviado. El docente indicó que se trata de más de un acto en el mismo momento referido a un inmueble, y que por eso existe el inciso de simultaneidad. % TODO: contenido incompleto o ambiguo en las notas originales (pasaje inaudible)

\subsection{Cuadros comparativos}

\begin{table}[H]
\centering
\caption{Hipoteca: las tres alternativas del comprador}
\small
\begin{tabularx}{\textwidth}{>{\raggedright\arraybackslash}p{2.6cm}>{\raggedright\arraybackslash}X>{\raggedright\arraybackslash}X>{\raggedright\arraybackslash}X}
\toprule
\textbf{Alternativa} & \textbf{En qué consiste} & \textbf{Responsabilidad del comprador} & \textbf{Conformidad del acreedor} \\
\midrule
1. Cancelar & El vendedor cancela su deuda al vender: acto uno, cancelo; acto dos, vendo libre. & Ninguna: recibe el inmueble libre de gravamen. & No es necesaria. \\
2. Reconocer y aceptar & El adquirente reconoce el gravamen; el deudor original sigue pagando las cuotas. & Solo ante una ejecución hipotecaria y hasta el valor del bien (inmueble 100, deuda 120: no debe 20). & No se libera al deudor original. \\
3. Asumir la deuda & El comprador ocupa el lugar del deudor inicial. & Todo su patrimonio, prenda común del acreedor (inmueble 100, deuda 120: responde por 120). & Solo si hay delegación pasiva (exonera al deudor original); si no asiste, el acreedor tiene dos patrimonios para ejecutar. \\
\bottomrule
\end{tabularx}
\end{table}

\begin{table}[H]
\centering
\caption{Derechos reales de disfrute y de garantía}
\small
\begin{tabularx}{\textwidth}{>{\raggedright\arraybackslash}p{2.6cm}>{\raggedright\arraybackslash}X>{\raggedright\arraybackslash}X}
\toprule
\textbf{Aspecto} & \textbf{Derechos reales de disfrute} & \textbf{Derechos reales de garantía} \\
\midrule
Ejemplos & Usufructo, uso, habitación y servidumbre (de paso, administrativas). & Hipoteca (de un banco o del propio vendedor por saldo de precio). \\
¿Se puede vender? & Sí (art.\ 1009): se vende la nuda propiedad o el derecho de usufructo. & Sí (art.\ 1009), con alguna de las tres alternativas. \\
Qué debe hacer el adquirente & Notificar a la otra parte; no se requiere conformidad. En la servidumbre, reconocerla y aceptarla. & Recibir el inmueble libre, reconocer y aceptar la hipoteca, o asumir la deuda. \\
Rol del tercero & El usufructuario tiene un derecho de menor jerarquía: solo debe estar notificado. & El acreedor presta conformidad solo para exonerar al deudor original. \\
\bottomrule
\end{tabularx}
\end{table}

\section*{Cuadro Cornell: derechos de terceros}
\addcontentsline{toc}{section}{Cuadro Cornell: derechos de terceros}

\begin{longtable}{>{\raggedright\arraybackslash}p{0.27\textwidth}>{\raggedright\arraybackslash}p{0.65\textwidth}}
\toprule
\textbf{Preguntas / palabras clave} & \textbf{Notas / respuestas} \\
\midrule
\endfirsthead
\toprule
\textbf{Preguntas / palabras clave} & \textbf{Notas / respuestas} \\
\midrule
\endhead
\textbf{¿Se puede vender un inmueble litigioso, gravado o con medidas cautelares? (art.\ 1009)} &
Sí: pueden ser objeto de los contratos sin perjuicio de los derechos de terceros, y quien contrata de mala fe como si estuvieran libres debe reparar los daños. «Gravado» implica un gravamen voluntario; la medida cautelar la impone un juez. Se distinguen tres situaciones: derechos reales de disfrute, derechos reales de garantía y medidas cautelares. \\[0.6em]
\textbf{Derechos reales de disfrute (usufructo, uso, habitación): ¿qué hay que hacer?} &
Se vende la nuda propiedad (dominio desmembrado). No se requiere conformidad del usufructuario, solo notificarlo, porque el nudo propietario dispone libremente y el usufructuario tiene un derecho de menor jerarquía. Lo mismo si el usufructuario transmite su derecho (el nuevo debe reconocer al nudo propietario y dar garantía suficiente). Al fallecer un usufructuario vitalicio, el nudo propietario adquiere el dominio pleno, salvo derecho de acrecer. \\[0.6em]
\textbf{Servidumbres} &
De paso o administrativas (por ejemplo, una cochera o unidad complementaria a la que se accede por otra unidad). El comprador la reconoce y la acepta; en la escritura se indica que del certificado surge la servidumbre y se inscribe como acto el reconocimiento, por ser un derecho real de disfrute. \\[0.6em]
\textbf{Hipoteca: ¿cuáles son las tres alternativas?} &
1) Cancelar: el vendedor cancela y se vende libre. 2) Reconocer y aceptar: el adquirente responde solo con el inmueble, hasta el límite del bien (inmueble 100, deuda 120: no debe 20) y el deudor original sigue pagando. 3) Asumir la deuda: el comprador ocupa el lugar del deudor y responde con todo su patrimonio (responde por 120). \\[0.6em]
\textbf{¿Cuándo debe concurrir el acreedor hipotecario?} &
Solo cuando hay delegación pasiva de deuda, es decir, cuando exonera al deudor original. Si no asiste y el comprador asumió la deuda, el acreedor tiene dos patrimonios para ejecutar. En los derechos de disfrute basta la notificación; en los de garantía, la conformidad solo para exonerar. \\
\midrule
\multicolumn{2}{p{0.94\textwidth}}{\textbf{Resumen:} frente a derechos reales de disfrute basta notificar al titular del otro derecho, sin conformidad. Frente a la hipoteca hay tres caminos: cancelar, reconocer y aceptar, o asumir la deuda; el acreedor interviene solo para exonerar al deudor original.} \\
\bottomrule
\end{longtable}

% =====================================================
\chapter{Medidas cautelares sobre el inmueble}
% =====================================================

Las medidas cautelares no son voluntarias: las ordena un juez.

\section{Anotación de litis}

La anotación de litis es una medida cautelar que publicita en el Registro de la Propiedad Inmueble una situación litigiosa (un juicio) que recae sobre un inmueble.

El inmueble puede venderse aun existiendo una situación litigiosa. El conocimiento del juicio afecta la decisión del comprador: puede decidir comprar según el estado del expediente. Al estar publicitado en el Registro, se \textbf{derriba la buena fe} del comprador: si la sentencia le resulta desfavorable, no puede alegar que no sabía.

\begin{important}
La anotación de litis es una medida cautelar que publicita una situación litigiosa que recae sobre un inmueble, y su objetivo es derribar la buena fe del comprador. No busca otra cosa: el proceso sigue su curso. Se puede comprar, pero se asume el riesgo: si la sentencia resulta contraria al titular de dominio, el comprador no puede reclamar nada a nadie, ni al vendedor, porque contrató en esas condiciones.
\end{important}

Puede ocurrir que alguien compre un inmueble muy barato apostando a que la sentencia salga a su favor, lo que a lo mejor es un buen negocio.

Se mencionaron \textbf{dos excepciones} en las que un inmueble con anotación de litis no se puede vender: el juicio de prescripción adquisitiva (usucapión) y la expropiación (cuando el Estado, por razones de utilidad pública, expropia a un privado un inmueble, por ejemplo para una plaza o una autopista). En el aula se discutió si podría comprarse cuando el plazo de la prescripción aún no se cumplió y solo existe un planteamiento, cuestión que se vería en juicio. % TODO: contenido incompleto o ambiguo en las notas originales (pasajes inaudibles en la discusión sobre usucapión)

Si la expropiación se cae, el inmueble sigue siendo de dominio privado (se mencionó un caso de una propiedad para una obra de subte).

\section{Prohibición de innovar}

Es una medida cautelar que veda la posibilidad de realizar cualquier acto jurídico sobre el bien: no se puede modificar la situación física ni jurídica del inmueble. Si recae una prohibición de innovar, \textbf{no se puede vender} el inmueble. Constituye una primera restricción a la regla del art.\ 1009.

\section{Prohibición de contratar}

Es una medida cautelar que restringe al titular de dominio la posibilidad de celebrar cualquier acto jurídico con relación al inmueble. Sigue la misma regla: \textbf{no se puede vender} el inmueble.

\section{Embargo preventivo}

Interesa el embargo preventivo: busca que el bien no salga del patrimonio de la persona debido a que esta tiene una deuda. Se afecta el bien al pago de una deuda.

\subsection{Cómo se lee la matrícula y el informe de dominio}

El registro informa tres columnas: la \textbf{titularidad}, los \textbf{gravámenes y restricciones} (y afectaciones al dominio) y las \textbf{cancelaciones}. La matrícula contiene además la calle, la nomenclatura catastral y la descripción del inmueble.

\begin{itemize}
\item \textbf{Titularidad:} muestra el encadenamiento sucesivo de los titulares de dominio y el negocio causal (por ejemplo, compraventa o donación).
\item \textbf{Gravámenes, restricciones y limitaciones:} de aquí surgen los derechos reales de disfrute, los de garantía y las medidas cautelares, así como las afectaciones al dominio (por ejemplo, el régimen de protección a la vivienda).
\item \textbf{Cancelaciones:} si el crédito hipotecario ya está abonado, figura tachado y en la parte de cancelaciones aparece la cancelación de la hipoteca.
\end{itemize}

Así se lee un certificado o informe de dominio, que muestra la realidad jurídica del inmueble. Existe una diferencia entre informe y certificado, pero la matrícula es siempre la misma.

\subsection{Situaciones frente al embargo}

Si el informe muestra un embargo, informa el monto y los autos del expediente. Lo primero es \textbf{ir a ver el expediente judicial}, porque puede ocurrir que el deudor ya haya abonado la deuda y el juez no haya librado oficio al Registro. El embargo está vigente a menos que se inscriba en el Registro el oficio judicial que ordena su cancelación.

\begin{enumerate}
\item \textbf{Embargo cancelado:} se deja constancia en la escritura; el juez envía el oficio al Registro, se cancela el embargo y no hay más complejidad (igual que la primera situación de la hipoteca).
\item \textbf{Embargo con deuda impaga, reconocerlo y aceptarlo:} si se transmite el dominio con el embargo, el adquirente lo acepta y únicamente responde con el bien que adquiere.
\item \textbf{Embargo con deuda impaga, asumirlo:} igual que en la hipoteca, puede responderse con todo el patrimonio y colocarse en el lugar del deudor.
\end{enumerate}

\subsection{Embargo pagado pero no cancelado en el registro}

Si la deuda ya se pagó pero no está cancelado en el registro, hay que avisar al juez para que libre oficio y se presente el abogado que inició el juicio. Debe lograrse una resolución judicial que vaya al Registro para que en la matrícula figure la cancelación del embargo. La escritura se hace cuando ya está cancelado en el registro, y en ella se transcribe la resolución judicial que ordena el levantamiento.

Recomendación práctica: pedir primero un informe y, cuando en el expediente se ve que el embargo se levantó, recién pedir el certificado. Si se pide el certificado y aparece un embargo, el trámite de levantamiento agota los 15 días y será necesario pedir otro certificado. Ordenado el levantamiento hay dos opciones: (1) anotarlo en el registro por testimonio judicial, lo que demora de mínima dos meses hasta que se inscribe; o (2) transcribir la pieza judicial en la escritura.

\subsection{Inmueble afectado al régimen de protección de la vivienda}

Un inmueble afectado al régimen de protección de la vivienda \textbf{se puede vender}: primero se cancela la afectación (desafectación) y luego se vende; puede hacerse en el mismo acto (capítulo uno, desafectación; capítulo dos, compraventa), vendiéndolo libre.

Si el vendedor quiere comprar otro inmueble, puede pedir la \textbf{subrogación real} del régimen. Se toma como fecha inicial la afectación original, porque se traslada el beneficio y no se afecta de nuevo. Si se desafecta en la venta y se afecta nuevamente en la compra sin más, el plazo corre desde cero. Para el traslado debe hacerse constar que se trata de un traslado del beneficio, en virtud de la afectación originaria de fecha tal sobre el inmueble tal, desafectado el día tal, con reserva del beneficio. Si no se dice nada de ello, se cuenta desde la fecha actual.

\section{El plazo de cinco años y qué medidas permiten escriturar}

\begin{important}
\textbf{Primero, siempre, el plazo.} Al encontrar una anotación en la columna del medio, primero se verifica el plazo: \textbf{cinco años desde la presentación en el registro}, no desde la medida ni desde la orden judicial. Si ya pasaron los cinco años, se hace de cuenta que no hay nada. Si se está dentro del plazo, hay que analizar de qué medida se trata: las dos prohibiciones (no innovar y no contratar) no permiten escriturar; la anotación de litis y el embargo sí, pero el adquirente debe tomar posición (en la litis, declarar que la conoce; en el embargo, reconocer o asumir). Se cita como fundamento que el Registro tiene por finalidad publicitar para mantener la vigencia de los derechos contra terceros interesados y de buena fe.
\end{important}

\begin{table}[H]
\centering
\caption{Medidas cautelares sobre inmuebles: ¿se puede vender?}
\small
\begin{tabularx}{\textwidth}{>{\raggedright\arraybackslash}p{2.3cm}>{\raggedright\arraybackslash}X>{\raggedright\arraybackslash}p{1.8cm}>{\raggedright\arraybackslash}X}
\toprule
\textbf{Medida} & \textbf{Qué es} & \textbf{¿Se vende?} & \textbf{Qué debe hacer el adquirente / observaciones} \\
\midrule
Anotación de litis & Publicita en el Registro una situación litigiosa; busca derribar la buena fe del comprador. & Sí, con riesgo & Conocerla y asumir el riesgo: si la sentencia es contraria al titular, no puede reclamar nada ni al vendedor. Excepciones señaladas: usucapión y expropiación. \\
Prohibición de innovar & Veda realizar actos jurídicos y modificar la situación física o jurídica del inmueble. & No & Restricción a la regla del art.\ 1009. \\
Prohibición de contratar & Restringe al titular la posibilidad de celebrar cualquier acto jurídico sobre el inmueble. & No & Misma regla que la prohibición de innovar. \\
Embargo preventivo & Afecta el bien al pago de una deuda para que no salga del patrimonio del deudor. & Sí & Cancelado: constancia en la escritura. No cancelado: reconocer y aceptar (responde con el bien) o asumir (responde con todo su patrimonio). Pagado pero vigente: reconocer «al solo efecto registral». \\
\bottomrule
\end{tabularx}
\end{table}

\section{Embargo pagado pero vigente en el registro: «al solo efecto registral»}

Puede ocurrir que el embargo esté vigente a nivel de publicidad, pero la deuda ya se haya abonado (por ejemplo, se pagó en el año dos y se está en el año tres), y los vendedores no gestionen el levantamiento. Si el comprador quiere seguir adelante, puede comprar. El consejo profesional es que revise el expediente con su abogado y se cerciore de que la deuda fue abonada (hay recibo de la parte acreedora, pero nadie solicitó el levantamiento al juez y el expediente quedó inactivo, ya que los procesos siguen a instancia de parte). No hay riesgo jurídico, porque no hay deuda: el comprador quiere comprar un inmueble, no una deuda. La opción es \textbf{comprar reconociendo el embargo al solo efecto registral}.

\begin{important}
«Al solo efecto registral» es una expresión que no existe en ninguna ley. Se usa porque, si la escritura no dijera nada sobre el embargo, el Registro la observaría y la devolvería sin registrarla, ya que figura un embargo vigente y no se tiene el levantamiento para transcribir la pieza judicial. El comprador declara que reconoce el embargo solo a efectos de que el registrador sepa, nada más; pero debe revisar el expediente con su letrado y cerciorarse de que esté pago.
\end{important}

Usualmente se hace declarar al vendedor, bajo juramento, que la totalidad de la deuda está abonada, como declaración adicional de la parte vendedora. Es común, porque en general los abogados cobran para obtener un levantamiento y tramitar el testimonio, y acreedor y deudor suelen olvidarse del tema una vez que arreglan. Lo mismo pasa con las hipotecas ya pagadas: muchas veces se espera que caduque el asiento por el paso del tiempo para no gastar en una escritura de cancelación, aunque si se quiere vender antes habrá que hacerla.

Frente a la consulta sobre quién declara que se pagó: en el expediente, quien tiene que recibir el dinero debe decir que ya lo recibió, y en función de eso el juez ordena el levantamiento. En cuanto a la escritura, la declaración del vendedor es un negocio que no pertenece al escribano, que puede ni saber de qué se trataba la deuda; por eso, ante un embargo vigente, hay que ir a ver el expediente.

\section{Conclusiones sobre inmuebles litigiosos, gravados o con medidas cautelares}

\begin{important}
Los inmuebles de dominio privado se pueden vender; los de dominio público, no. Con derechos reales puede celebrarse el contrato sin obstáculo, pero el adquirente debe hacer o decir algo: notificar al usufructuario, usuario o habitador, o, en los de garantía, reconocer la hipoteca o asumir la deuda. En las medidas cautelares hay más complejidad porque hay alguien que no está de acuerdo; si hay un expediente se trata de bienes litigiosos, que requieren mayor cuidado. Lo primero es ver el plazo de vigencia: pasados cinco años, se hace de cuenta que no existe nada. Dentro del plazo, hay que ver de qué medida se trata. Si se puede seguir adelante, el adquirente necesariamente debe hacer o decir algo, porque no se puede perjudicar a terceros y no se puede transmitir un derecho mejor y más extenso que el que se tiene: si el inmueble está embargado, no se lo puede transmitir libre.
\end{important}

\section*{Cuadro Cornell: medidas cautelares}
\addcontentsline{toc}{section}{Cuadro Cornell: medidas cautelares}

\begin{longtable}{>{\raggedright\arraybackslash}p{0.27\textwidth}>{\raggedright\arraybackslash}p{0.65\textwidth}}
\toprule
\textbf{Preguntas / palabras clave} & \textbf{Notas / respuestas} \\
\midrule
\endfirsthead
\toprule
\textbf{Preguntas / palabras clave} & \textbf{Notas / respuestas} \\
\midrule
\endhead
\textbf{¿Qué es la anotación de litis y qué efecto tiene?} &
Es una medida cautelar que publicita en el Registro una situación litigiosa; su objetivo es derribar la buena fe del comprador. Se puede comprar, pero se asume el riesgo: si la sentencia es contraria al titular de dominio, no se puede reclamar nada a nadie, ni al vendedor. Como excepciones en que no se puede vender se mencionaron el juicio de usucapión y la expropiación. \\[0.6em]
\textbf{Prohibición de innovar y prohibición de contratar} &
Impiden realizar actos jurídicos sobre el bien o modificar su situación física o jurídica. Si una de ellas está vigente, el inmueble no puede venderse: es una restricción a la regla del art.\ 1009. \\[0.6em]
\textbf{Embargo preventivo: ¿cómo se procede?} &
Se pide el informe, se lee la matrícula (titularidad, gravámenes y restricciones, cancelaciones) y se va al expediente judicial. Si está cancelado, el juez libra oficio al Registro y se deja constancia en la escritura. Si no está cancelado, se reconoce y acepta (responde con el bien) o se asume (responde con todo el patrimonio). Para levantarlo: testimonio judicial inscripto (mínimo dos meses) o transcripción de la resolución judicial en la escritura. \\[0.6em]
\textbf{Embargo pagado pero vigente: ¿qué es reconocerlo «al solo efecto registral»?} &
Si la deuda se pagó pero nadie pidió el levantamiento, el comprador reconoce el embargo al solo efecto registral para que el Registro inscriba la escritura. Es una fórmula usual, no prevista en ninguna ley. El comprador debe revisar el expediente con su letrado y suele hacerse declarar al vendedor, bajo juramento, que la deuda está abonada. \\[0.6em]
\textbf{¿Cuánto duran las medidas cautelares y qué se analiza primero?} &
Cinco años desde la presentación en el registro (no desde la orden judicial). Primero se fija el plazo: vencido, se hace de cuenta que no existe. Dentro del plazo, las prohibiciones de innovar y de contratar impiden escriturar; la anotación de litis y el embargo lo permiten, pero el adquirente debe tomar posición. \\[0.6em]
\textbf{Protección de la vivienda} &
El inmueble afectado se puede vender: primero se desafecta (acto uno, desafectación; acto dos, compraventa). Por subrogación real puede trasladarse el beneficio al nuevo inmueble: si se hace constar el traslado, se mantiene la fecha de la afectación originaria; si se desafecta y se afecta de nuevo sin decir nada, el plazo cuenta desde cero. \\
\midrule
\multicolumn{2}{p{0.94\textwidth}}{\textbf{Resumen:} las medidas cautelares exigen más cuidado porque hay una parte con intereses opuestos. Primero se verifica la vigencia de cinco años desde la presentación en el registro; las prohibiciones de innovar y de contratar impiden vender; la anotación de litis y el embargo lo permiten, el primero asumiendo el riesgo de perder la buena fe y el segundo reconociendo o asumiendo la deuda, o reconociéndolo «al solo efecto registral» cuando ya se pagó pero sigue inscripto.} \\
\bottomrule
\end{longtable}

% =====================================================
\chapter{Venta de cosas futuras e individualización del inmueble}
% =====================================================

La venta de cosas futuras es un tema complejo del que en esta clase solo se dejaron unas primeras ideas.

\section{Los inmuebles que se venden: lotes y unidades funcionales}

Jurídicamente, los inmuebles que se venden son básicamente de dos tipos: \textbf{unidades funcionales} o \textbf{lotes}. Que dentro del lote haya una casa o un baldío es lo mismo en términos jurídicos; solo varía el precio, el valor de mercado del bien. A efectos de la descripción, no se escribe la superficie cubierta del lote, sino el perímetro de lo que se vende. En el ámbito de la Capital se está más acostumbrado a la unidad funcional; en el resto del país no.

\section{¿Qué datos identifican jurídicamente una unidad funcional?}

Mediante una lluvia de ideas, los alumnos propusieron: nomenclatura, ubicación (calle, localidad), número de unidad funcional, balcón, frente o contrafrente, piso, superficie, entre calles, ambientes, medidas, dominio, reglamento y designación interna (por ejemplo, tercero A). El objetivo fue detectar cuáles datos son necesarios para identificar jurídicamente el bien y cuáles son adicionales, que pueden consignarse pero cuya ausencia no genera problema.

Datos que \textbf{no} identifican jurídicamente el bien: que esté amoblado (puede cambiar el aspecto económico, no la designación); la valuación (no es el objeto en sí, es otra variable); el balcón y los ambientes; el SUM (sigue siendo parte común); el frente o contrafrente.

\subsection{La designación interna}

No es obligatoria y no hace a la identificación jurídica, porque es un acuerdo de voluntades de los copropietarios que puede modificarse (por ejemplo, «las A pasan a ser las F»).

\subsection{El número de unidad funcional}

Define cuál es la unidad: no hay dos unidades funcionales dos ni dos unidades funcionales cinco. Es necesario para identificar el bien, pero por sí solo no alcanza, porque puede ser la unidad funcional cinco de la calle Corrientes o de la avenida Córdoba, de otro edificio. Hay que identificar dónde está asentado el bien.

\section{La nomenclatura catastral}

Con la nomenclatura catastral se identifica una parcela específica de una localidad. Es, por decirlo así, el «DNI» del inmueble. La otorga un organismo público distinto del Registro de la Propiedad Inmueble: el Registro se dedica a registrar ciertos actos jurídicos sobre ciertos bienes, mientras que este otro organismo se dedica al relevamiento físico de los inmuebles, que identifica no por calle y número sino por nomenclatura catastral. Según la jurisdicción se compone de más o menos datos. En la Ciudad de Buenos Aires tiene cuatro: \textbf{circunscripción, sección, manzana y parcela}, de lo más grande a lo más chico.

\begin{example}{Cómo funciona la nomenclatura}
La Ciudad de Buenos Aires está dividida en circunscripciones (el docente no recordó cuántas son; se mencionaron 19, 20 o 21). Supóngase la circunscripción cinco, que se divide en secciones: sección uno. La sección se divide en manzanas: manzana 10. Dentro de la manzana están las parcelas, cada una correspondiente a una casa o al lote de un edificio: parcela 33. Esto ubica específicamente un sector del territorio. % TODO: contenido incompleto o ambiguo en las notas originales (cantidad de circunscripciones)
\end{example}

En la provincia, además, hay otros datos: chacra, quinta y número del municipio.

\section{Descripción completa de la unidad funcional, calle y entre calles}

Para describir adecuadamente la unidad funcional se agregan datos particulares: el piso, la superficie y el dominio que le corresponde. La nomenclatura catastral pertenece a una jurisdicción determinada, por lo que se necesita saber la ciudad.

Calle, número de calle y entre calles se consignan por lo general, pero puede no haberlos: la calle puede no tener número, o encontrarse escrituras que venden un lote con frente a una calle sin nombre, sin número, entre otras dos calles sin nombre. Esto no preocupa, porque lo que identifica es la nomenclatura y, en el caso de las unidades funcionales, el número de unidad funcional. Los nombres de las calles pueden cambiar con el tiempo, por lo que no hacen a la identificación; si se tienen, se consignan porque ayudan a interpretar de qué se trata.

\begin{table}[H]
\centering
\caption{¿Qué identifica jurídicamente a una unidad funcional?}
\small
\begin{tabularx}{\textwidth}{>{\raggedright\arraybackslash}p{3.6cm}>{\raggedright\arraybackslash}p{2cm}>{\raggedright\arraybackslash}X}
\toprule
\textbf{Dato} & \textbf{¿Identifica?} & \textbf{Razón dada en clase} \\
\midrule
Número de unidad funcional & Sí & No hay dos unidades funcionales con el mismo número; por sí solo no alcanza. \\
Nomenclatura catastral & Sí & Identifica una parcela específica (circunscripción, sección, manzana y parcela): es el «DNI» del inmueble. \\
Jurisdicción (ciudad o municipio) & Sí & La nomenclatura catastral pertenece a una jurisdicción determinada. \\
Piso, superficie y dominio & Describe & Datos particulares que se agregan para describir adecuadamente la unidad. \\
Designación interna (ej.\ «3.º A») & No & Es un acuerdo de voluntades de los copropietarios y puede cambiar. \\
Calle, número y entre calles & No esencial & Los nombres pueden cambiar o no existir; se consignan si se tienen. \\
Amoblado, valuación, balcón, ambientes, SUM, frente o contrafrente & No & Cambian el valor o las características, no la identidad jurídica del bien. \\
\bottomrule
\end{tabularx}
\end{table}

\section*{Cuadro Cornell: individualización del inmueble}
\addcontentsline{toc}{section}{Cuadro Cornell: individualización del inmueble}

\begin{longtable}{>{\raggedright\arraybackslash}p{0.27\textwidth}>{\raggedright\arraybackslash}p{0.65\textwidth}}
\toprule
\textbf{Preguntas / palabras clave} & \textbf{Notas / respuestas} \\
\midrule
\endfirsthead
\toprule
\textbf{Preguntas / palabras clave} & \textbf{Notas / respuestas} \\
\midrule
\endhead
\textbf{Lote y unidad funcional} &
Jurídicamente solo hay lotes o unidades funcionales; que haya una casa o un baldío solo cambia el precio. En la descripción se consigna el perímetro del inmueble, no la superficie cubierta. \\[0.6em]
\textbf{¿Qué identifica jurídicamente a una unidad funcional?} &
Sí identifican: el número de unidad funcional (no hay dos iguales) y la nomenclatura catastral con su jurisdicción. Se agregan piso, superficie y dominio para describirla. No identifican: estar amoblada, la valuación, el balcón, los ambientes, el SUM, el frente o contrafrente, la designación interna (acuerdo de voluntades) ni el nombre de la calle. \\[0.6em]
\textbf{¿Qué es la nomenclatura catastral?} &
Es el «DNI» del inmueble. La otorga un organismo distinto del Registro, dedicado al relevamiento físico de los inmuebles. En la Ciudad de Buenos Aires se compone de circunscripción, sección, manzana y parcela, de lo más grande a lo más chico (ejemplo: circunscripción 5, sección 1, manzana 10, parcela 33). En la provincia se agregan, además, chacra, quinta y número de municipio. \\[0.6em]
\textbf{¿Por qué la calle y las entre calles no son esenciales?} &
Los nombres de las calles pueden cambiar o ni siquiera existir (calle sin nombre ni número). Se consignan si se tienen porque ayudan a interpretar, pero no hacen a la identificación: lo que identifica es la nomenclatura y, en las UF, el número de unidad funcional. \\
\midrule
\multicolumn{2}{p{0.94\textwidth}}{\textbf{Resumen:} un inmueble se identifica jurídicamente por el número de unidad funcional y la nomenclatura catastral (circunscripción, sección, manzana y parcela), mientras que la designación interna, el nombre de la calle y las características físicas o económicas no lo identifican.} \\
\bottomrule
\end{longtable}

% =====================================================
\chapter{Repaso general Cornell}
% =====================================================

\begin{longtable}{>{\raggedright\arraybackslash}p{0.27\textwidth}>{\raggedright\arraybackslash}p{0.65\textwidth}}
\toprule
\textbf{Preguntas / palabras clave} & \textbf{Notas / respuestas} \\
\midrule
\endfirsthead
\toprule
\textbf{Preguntas / palabras clave} & \textbf{Notas / respuestas} \\
\midrule
\endhead
\textbf{¿Cuándo puede venderse un inmueble según la clase?} &
Los inmuebles de dominio privado pueden venderse; los de dominio público no. Los bienes litigiosos, gravados o con medidas cautelares también pueden venderse (art.\ 1009), sin perjudicar a terceros y sin transmitir un derecho mejor ni más extenso que el que se tiene. \\[0.6em]
\textbf{¿Qué debe hacer el adquirente en cada caso?} &
Derechos de disfrute: notificar, sin conformidad (en la servidumbre, reconocerla). Hipoteca: cancelar, reconocer y aceptar, o asumir la deuda. Anotación de litis: conocerla y asumir el riesgo. Embargo: reconocerlo o asumirlo, o reconocerlo «al solo efecto registral» si ya se pagó. Prohibiciones de innovar y de contratar: no se puede vender. \\[0.6em]
\textbf{¿Qué se mira primero ante una medida cautelar?} &
El plazo de cinco años desde la presentación en el registro; vencido, se hace de cuenta que no existe. \\
\midrule
\multicolumn{2}{p{0.94\textwidth}}{\textbf{Síntesis final:} un inmueble puede venderse siempre que su objeto sea lícito, posible, determinado o determinable, valuable económicamente y existente o de existencia posible (arts.\ 1129, 279, 1003, 1004 y 16), con la excepción de los bienes del dominio público. El art.\ 1009 fija una regla permisiva para los bienes litigiosos, gravados o con medidas cautelares, y por eso en cada caso el adquirente debe hacer o decir algo. Frente a los derechos de disfrute basta notificar; frente a la hipoteca hay tres caminos y el acreedor interviene solo para exonerar al deudor original; las medidas cautelares exigen verificar primero su vigencia de cinco años, y las prohibiciones de innovar y de contratar impiden vender. La clase cierra con el inicio de la venta de cosas futuras y la individualización: el número de unidad funcional y la nomenclatura catastral identifican jurídicamente el inmueble; la designación interna, el nombre de la calle y las características físicas o económicas no.} \\
\bottomrule
\end{longtable}

\end{document}
